\section{Discussion}\label{sec:discussion}

\paragraph{Why a computed belief}
The comparison of Evidence$\rightarrow$LLM, Belief$\rightarrow$LLM and \core{} separates three ingredients: good evidence, a belief
over hypotheses, and who decides. Calibrated evidence alone lets a 7B model localise as well as the statistical core, but not type
the fault (FI 20.4 versus 43.7): choosing among 28 types requires weighing many weak signals against each other and against their
background rates, which a likelihood does consistently and a language model does not. Once the posterior is given, the same model
adopts it almost always. The belief is therefore the part that carries the diagnosis, and it can be computed. This complements
belief-state agents \citep{luo2026pos,luo2026graphofstates}: their beliefs are written by the model at a cost of millions of tokens,
ours is computed from calibrated likelihoods in seconds, and an LLM can still read it.

\paragraph{Where an LLM helps}
Bounding the LLM did not make it a better decision maker: when the evidence does not separate two types, the model chooses the more
specific one and loses more often than it wins, at 7B and at 120B alike. Scale changes the narrative instead---the larger model's
chain cites the right values at the right steps almost as often as the computed chain. The practical design that follows is a
separation of duties: the decision and its evidence are computed and auditable, and the language model explains them. The results
also suggest how a later version could let an LLM contribute to the decision without the damage observed here: as a sensor whose
answers are calibrated like any other categorical observation (against the incident's own background rate), or as a typing step
that must defer to the posterior unless it quotes evidence that discriminates between the top types. Because these ideas were
formed after seeing the RCA100 results, they must be evaluated on data not used for their design.

\paragraph{Failure modes ruled out by construction}
Several reasoning failures that Lu et al.\ \citep{lu2026beyond} observe in LLM agents cannot occur in \gama{} by construction
(\cref{tab:failuremodes}): every signal is tested against a baseline, every entity's resource signals are computed, timing enters as
onset precedence, caller and callee roles are fixed by the trace structure and the load direction, and the whole hypothesis space is
scored at once rather than anchored on the first idea. The remaining errors are mostly limits of the evidence---node faults without a
distinctive node signal, database faults without database metrics---rather than of the reasoning.

\paragraph{Threats to validity}
\emph{Internal.} RCA100 was the development data; although each change was pre-registered and the final modules were developed on
half of the cases, the 103-case results are optimistic, and we make no held-out claim on them. The fingerprint tables were written
from the published fault definitions, but their additions were informed by development errors. \emph{Construct.} Our scorer
re-implements the published RCA100 protocol; the retrieval-based half of the headline RP is independent of the decision and
favours systems that return evidence items, which is why every comparison is also reported with the step-anchored variant. The
Final weights of the benchmark (0.4 / 0.3 / 0.3) are taken as given. \emph{Statistical.} With 103 cases, differences of about three
points are at the edge of detectability; we report intervals and exact case counts with every test. API decoding is not
deterministic; we report the single pre-registered run of each arm. \emph{External.} RCA100 covers one application on one platform;
the AIOps2025 transfer (\cref{sec:aiops}) tests a second system with different fault types and pod-level faults. Two backbones do not
span the space of LLMs; a stronger reasoning model or a different interface may behave differently.

\paragraph{Ethics and data}
No personal data are involved. Both benchmarks are used under their non-commercial licences and are not redistributed. The API runs
used a provider setting without data retention; prompts contain only benchmark telemetry.
